\section{Introduction}

Automated-driving systems are increasingly assessed through public-road collision and disengagement reports. In the United States, the National Highway Traffic Safety Administration requires selected manufacturers and operators to report certain crashes involving automated driving systems or Level 2 advanced driver assistance systems, while the California Department of Motor Vehicles publishes autonomous-vehicle collision and disengagement reports from public-road testing \citep{nhtsa2021sgo,californiadmv_collision,californiadmv_disengagement}. These reports are valuable for accident analysis because they describe real-world interactions among automated vehicles, road geometry, traffic participants, safety drivers, and terminal outcomes such as contact, collision, or manual intervention. They also expose a distinctive evidence problem: the external event may be visible while the driver-side control-and-feedback pathway remains only partially documented. A typical report may state that an automated vehicle was operating in autonomous mode at an intersection, that another vehicle approached or crossed its path, and that a collision or disengagement occurred. The same report may omit the driver-visible mode display, takeover demand, time budget, driver gaze, driver acknowledgement, manual-control trace, perception confidence, planner confidence, or exact HMI timing. These omissions make it unsafe to treat the report as a complete causal record of what the driver knew, how the driver updated that knowledge, or why the event occurred.

Existing automated-vehicle incident studies have shown that public reports can support useful crash-pattern, disengagement, and narrative analyses. Yet these analyses often remain at the level of event type, crash partner, road context, vehicle movement, injury outcome, or latent narrative theme. Human-factors studies, in contrast, show that driver situation awareness, automation-mode understanding, trust calibration, time budget, workload, and HMI feedback can shape takeover and intervention quality. The difficulty is that the measurements that make those studies strong, such as gaze, response time, simulator logs, controlled HMI manipulations, or driver-state measurements, are usually unavailable in public incident reports. System-theoretic safety analysis provides a language for control actions, feedback, process models, and mode confusion, but applying these concepts to sparse public incident text is not straightforward. A direct large language model (LLM), or a human analyst working without an explicit evidence boundary, can easily collapse from ``collision occurred'' to ``the driver failed to take over'', even when the report does not state that a takeover request was issued, perceived, ignored, delayed, or misunderstood.

The problem addressed in this paper is therefore not accident-cause reconstruction from sparse text. The problem is how to support driver-centered post-incident safety review when the report contains enough scene evidence to motivate analysis, but not enough HMI, driver-state, transition, or ADS-internal evidence to justify strong psychological or causal claims. We formulate this task as \emph{Evidence-Bounded Driver Process-Model Replay}: public automated-driving crash or disengagement text is transformed into a structured tabletop replay package that explicitly separates reported evidence, derived evidence, missing evidence, abductive pathway hypotheses, blocked claims, and minimum evidence requirements for stronger future review. In this formulation, the replay artifact is not a verdict about the true driver state. It is an auditable review object that asks what the driver's process model would need to contain, what the report actually supports, which process-model updates are unverifiable, which driver actions and unsafe-control-action pathways remain admissible or blocked, and which HMI or logging fields should be recorded in future incidents.

Three challenges make this task difficult. First, missing evidence must remain missing. Public reports often encode HMI and driver-state information as not reported, but non-reporting is not evidence that an HMI cue was absent or that the driver did not perceive it. Second, driver-centered unsafe control actions cannot be inferred directly from terminal outcomes. In STPA-HF terms, a driver action becomes unsafe only in a specific control context mediated by the driver's process model, update sources, control-action candidates, and feedback structure. Third, accident analysis requires a review artifact rather than a single label. Analysts need to inspect not only the top-ranked pathway, but also the competing pathways, blocked claims, evidence gaps, and reporting requirements that explain why a stronger conclusion cannot be drawn. These challenges are amplified for LLM-based analysis because language models are fluent at narrative completion but may overstate unsupported driver psychology or transform sparse reports into outcome-driven stories unless their reasoning is structurally constrained.

We propose an STPA-HF-constrained Evidence-Bounded Driver Process-Model Replay framework. The framework first builds provenance-aware evidence packets from incident text and maps source-supported facts into four driver process-model dimensions: controlled-process states (CPS), controlled-process behaviors (CPB), other-process states (OPS), and other-process behaviors (OPB). It then analyzes how those process-model variables could be formed or updated through reported or missing feedback sources, including HMI mode displays, takeover cues, capability-boundary information, external scene observations, and other-actor behavior cues. Candidate driver control actions are generated from the process-model and update structure, and unsafe-control-action pathways are derived forward from those actions rather than backward from the reported outcome. The outcome is used only as a compatibility check. A deterministic evidence gate and an LLM judge then rank admissible, weakly supported, and blocked pathways. The final output is a tabletop replay package containing evidence profiles, process-model nodes, update analysis, candidate actions, unsafe-control-action pathways, pathway rankings, blocked claims, replay questions, and minimum HMI/logging requirement candidates. A richer-evidence replay case study is reserved for the end of the experimental section to test whether the same replay structure becomes more specific when additional temporal, HMI, driver-action, or scene-evolution evidence reduces the gaps present in sparse text reports.

The paper makes three contributions. First, it defines Evidence-Bounded Driver Process-Model Replay as a post-incident safety-review task for automated-driving crash and disengagement reports, reframing sparse reports as evidence-limited inputs rather than complete causal records (Section~\ref{sec:method}). Second, it develops an STPA-HF-constrained LLM framework that operationalizes driver process-model variables, update sources, candidate driver actions, unsafe-control-action pathways, outcome compatibility, blocked claims, and replay-package generation (Section~\ref{sec:method}). Third, it evaluates the framework on public automated-driving incident text by testing replay-package generation quality, outcome-only overreach suppression against direct LLM and generic chain-of-thought baselines, expert-alignment of ranked replay pathways, minimum reporting/logging requirement extraction, and HMI/logging sensitivity through counterfactual injections (Section~\ref{sec:experiments}). The final experiment reserves a richer-evidence case study as an evidence-density stress test. This component does not replace the text-report setting; it examines how the proposed replay artifact changes when more complete temporal, driver-action, HMI, or scene-evolution evidence is available.
